\documentclass[11pt,a4paper]{article}
\usepackage{shared/luxcover}
\usepackage[utf8]{inputenc}
\usepackage[T1]{fontenc}
\usepackage{amsmath,amssymb,amsfonts,amsthm}
\usepackage{graphicx}
\usepackage{hyperref}
\usepackage{booktabs}
\usepackage{algorithm}
\usepackage{algpseudocode}
\usepackage{enumitem}
\usepackage{xcolor}

\hypersetup{colorlinks=true,linkcolor=black,citecolor=blue,urlcolor=blue}

\newtheorem{theorem}{Theorem}[section]
\newtheorem{lemma}[theorem]{Lemma}
\newtheorem{definition}{Definition}[section]
\newtheorem{remark}{Remark}[section]

\title{LP-305: Quantum Secure Assets\\[4pt]
\large A Post-Quantum Methodology for Protecting Digital Assets\\
on the Lux Network}

\author{
Zach Kelling\thanks{Corresponding author: zach@lux.network}\\
\textit{Lux Industries, Inc.}\\
\texttt{research@lux.network}
}

\date{Original: August 2022\quad|\quad Revised: March 2026}

\begin{document}

\luxcoverpage

\begin{abstract}
We present the cryptographic architecture underlying Lux Network's
post-quantum asset protection layer. The design combines lattice-based
threshold cryptography (TFHE over RLWE), post-quantum signature schemes
(ML-DSA, Corona), fully homomorphic encryption for privacy-preserving
computation, and multimodal biometric authentication---providing
comprehensive security against adversaries with access to fault-tolerant
quantum computers.

This paper consolidates and updates the 2022 ``Quantum Secure Assets''
internal design document. Where the original described aspirational
architecture, this revision maps each component to its production
implementation, cites the formal verification where available, and
identifies the remaining open problems.
\end{abstract}

\tableofcontents
\newpage

%% ============================================================
\section{Introduction}

\subsection{Background}

Shor's algorithm~\cite{shor1999} renders RSA, ECDSA, and all
discrete-logarithm-based signature schemes insecure on a sufficiently
powerful quantum computer. NIST's Post-Quantum Cryptography
standardization process~\cite{nist-pqc} selected ML-KEM (Kyber),
ML-DSA (Dilithium), and SLH-DSA (SPHINCS+) as the primary
replacements. Lux Network adopted these standards at the consensus
layer and extended them with ring-signature privacy (Corona) and
fully homomorphic encryption (TFHE) for computation over ciphertexts.

\subsection{Scope}

This paper covers five layers of Lux's post-quantum stack:

\begin{enumerate}[nosep]
  \item \textbf{Consensus signatures}: BLS-381 + ML-DSA-65 + Corona
        (Quasar protocol, LP-105).
  \item \textbf{Threshold key management}: TFHE key generation via
        Shamir LSSS over a 2048-bit safe prime, with Feldman VSS
        commitments.
  \item \textbf{Fully homomorphic encryption}: TFHE boolean-gate
        circuits for encrypted CRDT merge, encrypted search, and
        private smart contracts.
  \item \textbf{Non-custodial wallets}: 2-of-3 MPC (CGGMP21 for
        secp256k1, FROST for Ed25519) with HSM co-signing.
  \item \textbf{Regulatory disclosure}: viewing keys, threshold
        disclosure, and ZK attestation anchoring.
\end{enumerate}

\subsection{Evolution from the 2022 Design}

The 2022 internal document proposed Lamport one-time signatures for
vault protection and an OTP-based information-theoretic signing layer.
Both were superseded:

\begin{itemize}[nosep]
  \item Lamport signatures $\rightarrow$ ML-DSA-65 (FIPS~204): multi-use,
        smaller signatures (3.3\,KB vs 16\,KB per Lamport instance),
        NIST-standardized.
  \item Network-layer OTP signing $\rightarrow$ dropped. The throughput
        cost of per-transaction OTP generation was incompatible with
        sub-second finality. ML-DSA provides computational (not
        information-theoretic) PQ security, which is sufficient given
        current lattice-hardness estimates~\cite{nist-pqc-report}.
\end{itemize}

%% ============================================================
\section{Lattice-Based Threshold Cryptography}

\subsection{Mathematical Foundation}

Security rests on the hardness of lattice problems. We define the
foundational problems before specializing to RLWE.

\begin{definition}[Lattice]
A lattice $L$ is the set of all integral combinations of $n$ linearly
independent vectors $\mathbf{b}_1, \mathbf{b}_2, \ldots, \mathbf{b}_n
\in \mathbb{R}^n$:
\[
L = \left\{ \sum_{i=1}^{n} x_i \mathbf{b}_i : x_i \in \mathbb{Z}
\right\}
\]
The matrix $B = [\mathbf{b}_1 | \cdots | \mathbf{b}_n]$ is a basis
of $L$.
\end{definition}

\begin{definition}[Shortest Vector Problem (SVP)]
Given a lattice $L$ and a norm $\|\cdot\|$, find a non-zero vector
$\mathbf{v} \in L$ that minimizes the norm:
\[
\mathbf{v} \in L \setminus \{\mathbf{0}\} \quad\text{such that}\quad
\|\mathbf{v}\| = \lambda_1(L) \triangleq
\min_{\mathbf{0} \neq \mathbf{u} \in L} \|\mathbf{u}\|
\]
The $\gamma$-approximate SVP ($\gamma$-SVP) relaxes this to finding
$\mathbf{v}$ with $\|\mathbf{v}\| \leq \gamma \cdot \lambda_1(L)$.
\end{definition}

\begin{definition}[Closest Vector Problem (CVP)]
Given a lattice $L$, a norm $\|\cdot\|$, and a target
$\mathbf{t} \in \mathbb{R}^n$, find a lattice vector closest
to $\mathbf{t}$:
\[
\mathbf{v} \in L \quad\text{such that}\quad
\|\mathbf{v} - \mathbf{t}\| =
\min_{\mathbf{u} \in L} \|\mathbf{u} - \mathbf{t}\|
\]
\end{definition}

Both SVP and CVP are NP-hard under randomized
reductions~\cite{ajtai1996}. No polynomial-time quantum algorithm is
known for either problem, even approximately with sub-exponential
approximation factors.

\begin{definition}[Learning With Errors (LWE)~\cite{regev2005}]
For dimension $n$, modulus $q$, and error distribution $\chi$ over
$\mathbb{Z}_q$, the LWE problem is: given a secret
$\mathbf{s} \in \mathbb{Z}_q^n$ and samples
$(\mathbf{a}_i, b_i = \langle \mathbf{a}_i, \mathbf{s} \rangle
+ e_i \bmod q)$ where $\mathbf{a}_i \xleftarrow{\$} \mathbb{Z}_q^n$
and $e_i \leftarrow \chi$, recover $\mathbf{s}$.

Equivalently, given matrix $A \in \mathbb{Z}_q^{m \times n}$ and
vector $\mathbf{b} \in \mathbb{Z}_q^m$:
\[
\text{Find } \mathbf{s} \in \mathbb{Z}_q^n
\text{ such that }
\|\mathbf{b} - A\mathbf{s}\| < \epsilon
\]
where $\epsilon$ bounds the error norm. Regev~\cite{regev2005} showed
that solving LWE is as hard as solving worst-case lattice problems
($\widetilde{O}(n/\alpha)$-approximate SVP) with quantum reductions.
\end{definition}

\subsubsection{Ring LWE}

The Ring LWE (RLWE) variant provides efficiency by working over
polynomial rings. Let $R_q = \mathbb{Z}_q[X]/(X^n + 1)$ be a
cyclotomic polynomial ring. The decisional RLWE problem asks to
distinguish $(a_i, a_i \cdot s + e_i)$ from uniform, where
$s \in R_q$ is secret and $e_i$ are drawn from a discrete Gaussian
$\chi_\sigma$.

\begin{definition}[RLWE Hardness Assumption]
For security parameter $\lambda$, ring dimension $n$, modulus $q$,
and error distribution $\chi_\sigma$, no probabilistic
polynomial-time (or bounded quantum polynomial-time) algorithm can
distinguish RLWE samples from uniform with advantage better than
$\mathrm{negl}(\lambda)$.
\end{definition}

RLWE reduces to ideal lattice problems in $R_q$, inheriting the
worst-case hardness of SVP over ideal lattices~\cite{lyubashevsky2013}.

\subsection{Implementation: \texttt{luxfi/fhe}}

Lux's TFHE library (\texttt{github.com/luxfi/fhe}) implements:

\begin{itemize}[nosep]
  \item \textbf{Parameter sets}: PN10QP27 ($N{=}1024$, $Q{=}134215681$,
        $\sim$128-bit classical), PN11QP54 ($N{=}2048$, higher precision),
        and OpenFHE-compatible STD128Q (post-quantum 128-bit).
  \item \textbf{Boolean gates}: AND, OR, XOR, NOT, NAND, NOR, XNOR,
        MUX, MAJORITY---each with programmable bootstrapping.
  \item \textbf{Integer arithmetic}: FheUint8/16/32/64 with
        Add, Sub, Mul (schoolbook), comparison chains (Lt, Gt, Eq).
  \item \textbf{Threshold key management}: Shamir secret sharing over
        the 2048-bit RFC~3526 Group~14 safe prime, with Feldman VSS
        commitments in the order-$q$ subgroup ($g = 2^2 \bmod p$).
  \item \textbf{Proactive reshare}: additive sub-sharing via Lagrange
        coefficients---no party materializes the full secret during
        reshare.
\end{itemize}

\subsection{Threshold FHE Protocol}

Key generation follows a $t$-of-$n$ Shamir scheme. Each party $i$
receives a share $s_i = f(i)$ where $f$ is a random degree-$(t{-}1)$
polynomial with $f(0) = \mathit{sk}$. Reconstruction uses Lagrange
interpolation:

\[
\mathit{sk} = \sum_{i \in S} s_i \cdot \lambda_i(0),
\qquad
\lambda_i(x) = \prod_{j \in S \setminus \{i\}} \frac{x - j}{i - j}
\]

Reshare without secret materialization: each old party $i$ picks a
random polynomial $g_i$ of degree $(t'{-}1)$ with constant term
$\lambda_i \cdot s_i$. New share $s'_j = \sum_i g_i(j)$. Since
$\sum_i g_i(0) = \sum_i \lambda_i \cdot s_i = \mathit{sk}$, the new
shares reconstruct the same secret. No single party ever holds
$\mathit{sk}$ in cleartext.

%% ============================================================
\section{Fully Homomorphic Encryption}

\subsection{Homomorphic Property}

The defining property of FHE is that it supports both addition and
multiplication on ciphertexts. Formally, an FHE scheme
$(\mathsf{KeyGen}, \mathsf{Enc}, \mathsf{Dec}, \mathsf{Eval})$
satisfies~\cite{gentry2009}:

\begin{theorem}[Homomorphic Correctness]
For any circuit $C$ of depth $d$, keys $(pk, sk) \leftarrow
\mathsf{KeyGen}(1^\lambda)$, and plaintexts $x, y$:
\begin{align}
\mathsf{Dec}(sk, \mathsf{Eval}(pk, +, \mathsf{Enc}(pk,x),
  \mathsf{Enc}(pk,y))) &= x + y \\
\mathsf{Dec}(sk, \mathsf{Eval}(pk, \times, \mathsf{Enc}(pk,x),
  \mathsf{Enc}(pk,y))) &= x \cdot y
\end{align}
provided the accumulated noise does not exceed the modulus bound.
\end{theorem}

The noise growth per operation is the fundamental obstacle. For
LWE-based schemes~\cite{gsw2013}, multiplication increases noise
quadratically. Gentry's bootstrapping technique~\cite{gentry2009}
resets noise by homomorphically evaluating the decryption circuit,
enabling unbounded computation depth at the cost of one bootstrapping
operation per gate.

TFHE~\cite{tfhe2016} achieves gate-by-gate bootstrapping: each
boolean gate evaluation simultaneously computes the gate and refreshes
the ciphertext noise in a single operation, avoiding noise
accumulation entirely.

\subsection{TFHE Gate-by-Gate Bootstrapping}

Each boolean gate $G \in \{\text{AND}, \text{OR}, \text{XOR}, \ldots\}$
is evaluated via a single bootstrapping operation that simultaneously
computes the gate and refreshes the noise:

\[
\text{Bootstrap}(c, G) = \text{BlindRotate}(c) \cdot G_{\text{LUT}}
\]

where $G_{\text{LUT}}$ encodes the gate's truth table as a test
polynomial in $R_q$.

\subsection{Encrypted CRDT Merge}

CRDT merge under FHE enables conflict-free replication without
decryption. The LWW-Register merge circuit compares encrypted
timestamps MSB-to-LSB:

\begin{algorithmic}[1]
\State $\mathit{bGtA} \gets \bot$; $\mathit{eqSoFar} \gets \top$
\For{$i = \text{MSB}$ \textbf{downto} $0$}
  \State $\mathit{bitGt} \gets \text{ANDNY}(a_i, b_i)$
  \State $\mathit{bitEq} \gets \text{XNOR}(a_i, b_i)$
  \State $\mathit{bGtA} \gets \text{OR}(\mathit{bGtA},
         \text{AND}(\mathit{eqSoFar}, \mathit{bitGt}))$
  \State $\mathit{eqSoFar} \gets \text{AND}(\mathit{eqSoFar},
         \mathit{bitEq})$
\EndFor
\For{each bit $i$}
  \State $\mathit{result}_i \gets \text{MUX}(\mathit{bGtA}, b_i, a_i)$
\EndFor
\end{algorithmic}

Gate count: $O(\text{bitsTS})$ comparisons + $O(\text{bitsVal} +
\text{bitsTS})$ MUX selections. At PN10QP27 with 8-bit timestamps
and 8-bit values: $\sim$40 bootstraps $\approx$ 4 seconds.

\subsection{Performance Budget}

\begin{center}
\begin{tabular}{lrrr}
\toprule
\textbf{Operation} & \textbf{Gates} & \textbf{Time (PN10QP27)} & \textbf{Ciphertext} \\
\midrule
LWW Merge (8+8 bit) & 40 & 4\,s & 136\,KB \\
GCounter Max (32 bit/node) & 128 & 13\,s & 544\,KB \\
OR-Set Merge & 0 (tag union) & $<$1\,ms & tags only \\
\bottomrule
\end{tabular}
\end{center}

%% ============================================================
\section{Non-Custodial Wallet Architecture}

Every user receives a 2-of-3 MPC wallet on the Lux EVM:

\begin{enumerate}[nosep]
  \item \textbf{User device} --- key share derived from WebAuthn PRF.
  \item \textbf{ATS co-signer} --- auto-signs after compliance checks.
  \item \textbf{Backup (cold)} --- Shamir recovery via
        \texttt{luxfi/age} (X25519 + ML-KEM-768 X-Wing hybrid).
\end{enumerate}

Protocols: CGGMP21~\cite{cggmp21} for ECDSA (secp256k1), FROST for
EdDSA (Ed25519). HSM co-signing for settlement intents
(AWS KMS, GCP, Zymbit).

%% ============================================================
\section{Regulatory Disclosure}

Three composable modes, implemented in \texttt{Disclosure.sol}:

\begin{enumerate}[nosep]
  \item \textbf{Viewing keys} (Zcash-style): owner registers a
        per-viewer wrapped decryption key. Revocable on-chain; actual
        key revocation requires off-chain rekey.
  \item \textbf{Threshold disclosure} ($M$-of-$N$, regulator is one
        party): uses the existing \texttt{lsss.go} Shamir infrastructure.
        On-chain share submission with per-party dedup. Combination
        happens off-chain.
  \item \textbf{ZK attestation anchor}: accepts opaque proof bytes
        with a \texttt{claimType} tag. On-chain verifier registry
        (owner-gated) routes to Groth16/PlonK verifier contracts when
        deployed.
\end{enumerate}

%% ============================================================
\section{Open Problems}

\begin{enumerate}[nosep]
  \item \textbf{Deterministic keygen from seed}: resolved in
        \texttt{luxfi/lattice} v8 (seeded-PRNG \texttt{KeyGenerator}).
        Consensus validators now derive identical FHE keys from a shared
        seed, enabling threshold reshare without trusted setup.
  \item \textbf{FHE throughput}: $\sim$0.24 ops/sec at production
        parameters (PN10QP27). GPU acceleration via NTT (LP-203) is the
        primary mitigation path for high-frequency use cases.
  \item \textbf{Biometric + FHE fusion}: encrypting biometric templates
        under TFHE so servers can \emph{match} without seeing the
        template remains an active engineering direction.
\end{enumerate}

%% ============================================================
\section{Conclusion}

The design combines lattice-based threshold cryptography, TFHE-backed
homomorphic CRDT merge, non-custodial MPC wallets (CGGMP21 + FROST),
and three-mode regulatory disclosure. Every component is anchored in
Lux production code --- \texttt{luxfi/fhe}, \texttt{luxfi/lattice},
\texttt{luxfi/mpc}, \texttt{luxfi/age}, \texttt{luxfi/standard} --- with
formal verification artefacts in \texttt{luxfi/proofs}.

\bibliographystyle{plain}
\begin{thebibliography}{99}

\bibitem{shor1999}
P.~W.~Shor, ``Polynomial-time algorithms for prime factorization and
discrete logarithms on a quantum computer,'' \textit{SIAM Review},
vol.~41, no.~2, pp.~303--332, 1999.

\bibitem{nist-pqc}
NIST, ``Post-Quantum Cryptography Standardization,''
\url{https://csrc.nist.gov/projects/post-quantum-cryptography}, 2024.

\bibitem{nist-pqc-report}
NIST, ``Status Report on the Third Round of the NIST Post-Quantum
Cryptography Standardization Process,'' NISTIR~8413, 2022.

\bibitem{cggmp21}
R.~Canetti, R.~Gennaro, S.~Goldfeder, N.~Makriyannis, and U.~Peled,
``UC Non-Interactive, Proactive, Threshold ECDSA with Identifiable
Aborts,'' IACR ePrint 2021/060, 2021.

\bibitem{regev2005}
O.~Regev, ``On lattices, learning with errors, random linear codes,
and cryptography,'' \textit{STOC}, pp.~84--93, 2005.

\bibitem{gentry2009}
C.~Gentry, ``A fully homomorphic encryption scheme,'' Ph.D.\
dissertation, Stanford University, 2009.

\bibitem{lamport1979}
L.~Lamport, ``Constructing digital signatures from a one-way
function,'' Technical Report SRI-CSL-98, 1979.

\bibitem{rfc3526}
T.~Kivinen and M.~Kojo, ``More Modular Exponential (MODP)
Diffie-Hellman groups for Internet Key Exchange (IKE),'' RFC~3526,
2003.

\bibitem{ajtai1996}
M.~Ajtai, ``Generating hard instances of lattice problems (extended
abstract),'' \textit{STOC}, pp.~99--108, 1996.

\bibitem{micciancio-regev}
D.~Micciancio and O.~Regev, ``Lattice-based cryptography,'' in
\textit{Post-Quantum Cryptography}, Springer, pp.~147--191, 2009.

\bibitem{lpr2013}
V.~Lyubashevsky, C.~Peikert, and O.~Regev, ``On ideal lattices and
learning with errors over rings,'' \textit{Journal of the ACM},
vol.~60, no.~6, 2013.

\bibitem{gsw2013}
C.~Gentry, A.~Sahai, and B.~Waters, ``Homomorphic encryption from
learning with errors: Conceptually-simpler, asymptotically-faster,
attribute-based,'' \textit{CRYPTO}, 2013.

\bibitem{sphincs-plus}
D.~J.~Bernstein, A.~H\"{u}lsing, S.~K\"{o}lbl, R.~Niederhagen,
J.~Rijneveld, and P.~Schwabe, ``The SPHINCS+ signature framework,''
\textit{ACM CCS}, 2019.

\bibitem{tfhe2016}
I.~Chillotti, N.~Gama, M.~Georgieva, and M.~Izabach\`{e}ne, ``Faster
Fully Homomorphic Encryption: Bootstrapping in less than 0.1
Seconds,'' \textit{ASIACRYPT}, 2016.

\end{thebibliography}

\end{document}
